\begin{lemma}
\label{lemma-construct-strictly-perfect}
Let $X$ be a quasi-compact and quasi-separated scheme with the
resolution property.
Let $\mathcal{F}^\bullet$ be a bounded below complex of quasi-coherent
$\mathcal{O}_X$-modules representing a perfect object of
$D(\mathcal{O}_X)$. Then there exists a bounded complex
$\mathcal{E}^\bullet$ of finite locally free $\mathcal{O}_X$-modules
and a quasi-isomorphism $\mathcal{E}^\bullet \to \mathcal{F}^\bullet$.
\end{lemma}

\begin{proof}
Let $a, b \in \mathbf{Z}$ be integers such that $\mathcal{F}^\bullet$
has tor amplitude in $[a, b]$ and such that $\mathcal{F}^n = 0$ for
$n < a$. The existence of such a pair of integers
follows from Cohomology, Lemma \ref{cohomology-lemma-perfect}
and the fact that $X$ is quasi-compact.
If $b < a$, then $\mathcal{F}^\bullet$ is zero in the
derived category and the lemma holds.
We will prove by induction on
$b - a \geq 0$ that there exists a complex
$\mathcal{E}^a \to \ldots \to \mathcal{E}^b$
with $\mathcal{E}^i$ finite locally free
and a quasi-isomorphism $\mathcal{E}^\bullet \to \mathcal{F}^\bullet$.

\medskip\noindent
The base case is the case $b - a = 0$. In this case
$H^b(\mathcal{F}^\bullet) = H^a(\mathcal{F}^\bullet) =
\Ker(\mathcal{F}^a \to \mathcal{F}^{a + 1})$
is finite locally free. Namely, it is a finitely presented
$\mathcal{O}_X$-module of tor dimension $0$ and hence finite
locally free. See Cohomology, Lemmas \ref{cohomology-lemma-perfect} and
\ref{cohomology-lemma-finite-cohomology} and
Properties, Lemma \ref{properties-lemma-finite-locally-free}.
Thus we can take $\mathcal{E}^\bullet$ to be
$H^b(\mathcal{F}^\bullet)$ sitting in degree $b$.
The rest of the proof is dedicated to the induction step.

\medskip\noindent
Assume $b > a$. Observe that
$$
H^b(\mathcal{F}^\bullet) =
\Ker(\mathcal{F}^b \to \mathcal{F}^{b + 1})/
\Im(\mathcal{F}^{b - 1} \to \mathcal{F}^b)
$$
is a finite type quasi-coherent $\mathcal{O}_X$-module, see
Cohomology, Lemmas \ref{cohomology-lemma-perfect} and
\ref{cohomology-lemma-finite-cohomology}. Then we can find a finite type
quasi-coherent $\mathcal{O}_X$-module $\mathcal{F}$ and a map
$$
\mathcal{F} \longrightarrow
\Ker(\mathcal{F}^b \to \mathcal{F}^{b + 1})
$$
such that the composition with the projection onto
$H^b(\mathcal{F}^\bullet)$ is surjective.
Namely, we can write $\Ker(\mathcal{F}^b \to \mathcal{F}^{b + 1})$
as the filtered union of its finite type quasi-coherent submodules by
Properties, Lemma \ref{properties-lemma-quasi-coherent-colimit-finite-type}
and then one of these will do the job.
Next, we choose a finite locally free $\mathcal{O}_X$-module
$\mathcal{E}^b$ and a surjection $\mathcal{E}^b \to \mathcal{F}$ using
the resolution property of $X$. Consider the map of complexes
$$
\alpha : \mathcal{E}^b[-b] \to \mathcal{F}^\bullet
$$
and its cone $C(\alpha)^\bullet$, see
Derived Categories, Definition \ref{derived-definition-cone}.
We observe that $C(\alpha)^\bullet$ is nonzero only in degrees
$\geq a$, has tor amplitude in $[a, b]$ by
Cohomology, Lemma \ref{cohomology-lemma-cone-tor-amplitude},
and has $H^b(C(\alpha)^\bullet) = 0$ by construction.
Thus we actually find that $C(\alpha)^\bullet$ has tor amplitude
in $[a, b - 1]$. Hence the induction hypothesis applies to
$C(\alpha)^\bullet$ and we find a map of complexes
$$
(\mathcal{E}^a \to \ldots \to \mathcal{E}^{b - 1})
\longrightarrow
C(\alpha)^\bullet
$$
with properties as stated in the induction hypothesis. Unwinding
the definition of the cone this gives a commutative diagram
$$
\xymatrix{
\ldots \ar[r] &
\mathcal{E}^{b - 2} \ar[r] \ar[d] &
\mathcal{E}^{b - 1} \ar[r] \ar[d] &
0 \ar[r] \ar[d] &
\ldots \\
\ldots \ar[r] &
\mathcal{F}^{b - 2} \ar[r] &
\mathcal{F}^{b - 1} \oplus \mathcal{E}^b \ar[r] &
\mathcal{F}^b \ar[r] &
\ldots
}
$$
It is clear that we obtain a map of complexes
$(\mathcal{E}^a \to \ldots \to \mathcal{E}^b) \to \mathcal{F}^\bullet$.
We omit the verification that this map is a quasi-isomorphism.
\end{proof}